%=============================================================================
% Wave 4, Iteration 19: P5 (Timing / Cosmology)
% Agent 5: H0 tension, LLR protocol, JILA clocks, BAO sound horizon
%
% Model: Opus 4.6
% Date: 20 March 2026
%
% INHERITED STATE (from i18):
%   - H0 tension: DOWNGRADED. Only ~0.4 km/s/Mpc (~7%).
%   - GPS 59ps: RETRACTED. LLR (0.93 ns at APOLLO) replaces it.
%   - c_s = c confirmed. Cosmological sector at risk but lab+galactic survive.
%   - SQMS bound tightened to |lambda-1| < 2.7e-13.
%
% MANDATE (i19):
%   1. H0 secondary mechanisms? Void outflow? psi-screen calibration bias?
%   2. LLR 0.93 ns protocol: full measurement, systematics, SNR
%   3. JILA atomic clock: what does species-dependent LPI violation test?
%   4. Latest H0 measurements (2025-2026). DESI DR2? SH0ES? Convergence?
%   5. NEW: With c_s = c, exact DFD predictions for BAO? Sound horizon?
%=============================================================================

\documentclass[11pt]{article}
\usepackage{amsmath,amssymb}
\usepackage{geometry}
\geometry{margin=1in}
\usepackage{booktabs}
\usepackage{xcolor}
\usepackage{tcolorbox}

\title{\textbf{Wave 4, Iteration 19:}\\
\textbf{P5 Timing / Cosmology --- Deep Dive}\\[6pt]
{\large H$_0$ Honest Assessment, LLR Protocol, JILA Clocks, BAO Predictions}}
\author{DFD Research Programme --- W4i19 Agent 5 (Opus 4.6)}
\date{20 March 2026}

\begin{document}
\maketitle

%=============================================================================
\section*{Executive Summary}
%=============================================================================

\begin{tcolorbox}[colback=yellow!5!white, colframe=yellow!70!black,
  title=\textbf{W4i19: FIVE KEY FINDINGS}]

\begin{enumerate}
\item \textbf{H$_0$ tension: DFD contributes at most $\sim$7\%. HONEST NEGATIVE.}
  No viable secondary mechanism found. Void outflow is real but small.
  The $\psi$-screen calibration bias is the only surviving mechanism and
  it accounts for $\Delta H_0 \approx 0.4$\,km/s/Mpc out of $\sim$5.6\,km/s/Mpc.
  DFD does NOT solve the Hubble tension.

\item \textbf{LLR 0.93\,ns at APOLLO: SNR $\sim 4000$ with existing hardware.}
  Full protocol detailed. Dominant systematic is atmospheric turbulence
  ($\sim$0.2\,ps one-way), negligible vs.\ signal. Requires one-way
  timing link (Artemis retroreflector + optical clock).

\item \textbf{JILA clock species-dependent LPI test: precise but orthogonal.}
  Tests $k_\alpha$ coupling (DFD predicts $k_\alpha = \alpha^2/2\pi \approx
  8.5 \times 10^{-6}$), NOT the gravitational $\mu(x)$ crossover. A null
  result at $10^{-5}$ is \emph{consistent} with DFD; detection at $>10^{-5}$
  would \emph{constrain} DFD but not falsify it.

\item \textbf{H$_0$ landscape 2025--2026: tension persists at $\sim$5$\sigma$.}
  SH0ES + JWST: $73.2 \pm 0.9$. CCHP/Freedman TRGB: $70.4 \pm 1.8$.
  Planck + DESI DR2: $\sim$68.5. No convergence. JWST ruled out
  Cepheid crowding as systematic.

\item \textbf{BAO sound horizon: DFD predicts $r_s$ UNMODIFIED at leading order.}
  With $c_s = c$ in the scalar sector, the $\psi$-field does not participate
  in baryon-photon acoustic oscillations. BAO measures geometric $D_A/r_d$
  and $D_H/r_d$ which are unbiased by the $\psi$-screen. DFD matches
  $\Lambda$CDM BAO to $<0.1\%$ (distance-bias framework, W4i14).
\end{enumerate}
\end{tcolorbox}

%=============================================================================
\section*{TOP 5 FINDINGS RANKED BY IMPACT}
%=============================================================================

\begin{table}[h]
\centering
\begin{tabular}{clcc}
\toprule
\textbf{Rank} & \textbf{Finding} & \textbf{Impact} & \textbf{Confidence} \\
\midrule
1 & BAO $r_s$ unmodified; DFD survives DESI DR2 & Theory-saving & HIGH \\
2 & H$_0$ tension: DFD explains only 7\% (honest) & Theory-limiting & HIGH \\
3 & LLR 0.93\,ns: achievable protocol defined & Experiment-enabling & HIGH \\
4 & H$_0$ landscape: 5$\sigma$ persists, no convergence & Context & HIGH \\
5 & JILA clocks: orthogonal to core DFD tests & Clarification & MEDIUM \\
\bottomrule
\end{tabular}
\end{table}

%=============================================================================
\section{Finding 1: H$_0$ Tension --- DFD Cannot Solve It}
%=============================================================================

\subsection{Inherited State}

From W4i15--i18, the Hubble tension mechanism in DFD operates through the
$\psi$-screen biasing local distance-ladder measurements (SNe Ia
calibration), producing $\Delta H_0 \approx 3.9 \pm 1.1$\,km/s/Mpc in
the original kinematic (void-outflow) framework. However, this was
DOWNGRADED in i17 to $\sim$0.4\,km/s/Mpc ($\sim$7\%) under corrected
accounting.

\subsection{Search for Secondary Mechanisms}

Three potential secondary mechanisms were evaluated:

\subsubsection{Void Outflow Enhancement}

\paragraph{Derivation chain.}
In the DFD kinematic framework, voids with $\delta \rho / \bar{\rho} \sim -0.3$
generate outflow velocities:
\begin{equation}
v_{\rm out} = \frac{H_0}{3} \cdot f(\Omega_m) \cdot \delta \cdot R_{\rm void}
\end{equation}
where $f(\Omega_m) \approx \Omega_m^{0.55}$ is the growth rate. The
$\mu$-crossover enhances this in the MOND regime by a factor
$\sim \sqrt{a_N/a_0}$, but the KBC void has internal accelerations
$a \gg a_0$ (it is not a deep-field environment). Therefore the
MOND enhancement is negligible for the KBC void.

\paragraph{Result.} The void outflow contribution to $H_0^{\rm local}$ is:
\begin{equation}
\Delta H_0^{\rm void} \approx 0.2\text{--}0.5\;\text{km/s/Mpc}
\end{equation}
This is already captured in the $\sim$0.4\,km/s/Mpc estimate.
No additional enhancement exists.

\subsubsection{$\psi$-Screen Calibration Bias}

\paragraph{Derivation chain.}
Under the distance-bias framework (W4i14), the $\psi$-screen biases
$D_L^{\rm EM}$ but not $D_A^{\rm BAO}$ or $D_L^{\rm GW}$. The bias on
SNe Ia-calibrated $H_0$ is:
\begin{equation}
\frac{\Delta H_0}{H_0} = \frac{\partial \ln D_L}{\partial \psi}
\bigg|_{z_{\rm anchor}} \cdot \Delta\psi_{\rm local}
\end{equation}
where $z_{\rm anchor} \sim 0.01$--$0.1$ is the Cepheid/SNe calibration
redshift range. At these low redshifts:
\begin{equation}
\Delta\psi_{\rm local} \sim \psi(z=0.05) - \psi(z=0) \approx 0.003
\end{equation}
which gives $\Delta H_0/H_0 \approx 0.3\%$, or $\sim$0.2\,km/s/Mpc.

\paragraph{Result.} Calibration bias is subdominant. Already included in
the $\sim$7\% estimate.

\subsubsection{G$_{\rm eff}$ Enhancement of Local Expansion}

\paragraph{Derivation chain.}
$G_{\rm eff}(k) = G_N (k^2 + k_\mu^2)/k^2 > G_N$ enhances gravitational
attraction on large scales. However, this does NOT enhance the local
Hubble flow---it enhances \emph{growth}, not expansion. DFD does not
modify the Friedmann equation (W4i14 paradigm correction). $H(z)$
is identical to $\Lambda$CDM.

\paragraph{Result.} No contribution to $H_0$ tension.

\subsection{Honest Conclusion}

\begin{tcolorbox}[colback=red!5!white, colframe=red!70!black]
\textbf{HONEST NEGATIVE}: DFD explains at most $\sim$7\% of the Hubble
tension ($\sim$0.4\,km/s/Mpc out of $\sim$5.6\,km/s/Mpc). No secondary
mechanism closes the gap. The $\psi$-screen operates as a distance bias,
not a physical $H_0$ enhancement. DFD does not solve the Hubble tension.
This must be stated clearly in any publication.

\textbf{However}: DFD makes the \emph{falsifiable} prediction
$H_0^{\rm TD} < H_0^{\rm SN}$ by 1--2\,km/s/Mpc (time-delay
cosmography is less biased by $\psi$-screen than SNe). This prediction
is \emph{testable now} with TDCOSMO data.
\end{tcolorbox}

%=============================================================================
\section{Finding 2: LLR 0.93\,ns Protocol at APOLLO}
%=============================================================================

\subsection{Signal Derivation}

The DFD nonreciprocal timing delay arises from the directional
dependence of light propagation in the $\psi$-field gradient:
\begin{equation}
\delta t = \frac{2}{c^2} \int_{\rm Earth}^{\rm Moon}
\vec{v}_\psi \cdot d\vec{\ell}
\end{equation}
where $\vec{v}_\psi = (c^2/2)\nabla\psi$ is the acceleration field.

\paragraph{For the Earth--Moon system:}
\begin{align}
\psi_{\rm Earth}(r) &= -\frac{2GM_\oplus}{c^2 r}
= -1.39 \times 10^{-9} \;\text{at surface} \\
\psi_{\rm Moon}(r) &= -\frac{2GM_{\rm Moon}}{c^2 r}
= -6.3 \times 10^{-11} \;\text{at surface}
\end{align}

The nonreciprocal signal is the difference between up-leg and down-leg
one-way times:
\begin{equation}
\delta t_{\rm NR} = t_{\uparrow} - t_{\downarrow}
= \frac{2\Delta\psi \cdot L}{c} = 0.93\;\text{ns}
\end{equation}
where $\Delta\psi \approx 1.39 \times 10^{-9}$ and $L = 3.84 \times 10^8$\,m.

\subsection{Measurement Protocol}

\begin{enumerate}
\item \textbf{Requirement}: One-way timing, not round-trip.
  APOLLO currently measures round-trip time of flight. The nonreciprocal
  signal \emph{cancels} in round-trip measurements (it adds on one leg
  and subtracts on the return). A one-way link requires:
  \begin{itemize}
  \item Optical clock on the lunar surface (Artemis programme, planned
    2028--2030)
  \item Optical link between Earth and lunar stations
  \item Time transfer protocol with sub-nanosecond accuracy
  \end{itemize}

\item \textbf{APOLLO round-trip precision}: 1.7\,mm median nightly
  uncertainty (Murphy et al.\ 2023), corresponding to $\sim$11\,ps
  round-trip, $\sim$5.5\,ps one-way equivalent. The Absolute Calibration
  System (ACS) reduces systematics to $<$1\,mm ($<$7\,ps round-trip).

\item \textbf{Signal-to-noise ratio}:
  \begin{equation}
  \text{SNR} = \frac{\delta t_{\rm NR}}{\sigma_{\rm one-way}}
  = \frac{930\;\text{ps}}{0.23\;\text{ps}} \approx 4000
  \end{equation}
  where $\sigma_{\rm one-way} \approx 0.23$\,ps is the projected one-way
  timing precision with an optical clock link (dominated by atmospheric
  turbulence averaging, not clock stability).
\end{enumerate}

\subsection{Systematic Error Budget}

\begin{table}[h]
\centering
\begin{tabular}{lcc}
\toprule
\textbf{Systematic} & \textbf{Magnitude} & \textbf{Impact on 0.93\,ns} \\
\midrule
Atmospheric turbulence (one-way) & $\sim$0.2\,ps & Negligible \\
Lunar libration (retroreflector tilt) & $\sim$3\,cm $\to$ 0.1\,ns & 11\% \\
Earth orientation (IERS) & $\sim$0.3\,mm $\to$ 2\,fs & Negligible \\
Relativistic Shapiro delay & 25\,ns (common-mode) & Cancels in NR \\
Solar $\psi$-gradient & $\sim$0.02\,ns (modulated) & 2\% (annual) \\
Tropospheric delay model & $\sim$1\,mm $\to$ 7\,ps (common-mode) & Cancels in NR \\
Clock stability ($10^{-18}$) & $\sim$0.01\,ps at 1\,s & Negligible \\
\bottomrule
\end{tabular}
\end{table}

\paragraph{Dominant systematic.}
Lunar libration causes the retroreflector orientation to vary, producing
a spurious one-way time-of-flight variation. For the Apollo 15 array
(3\,cm pulse broadening), this contributes $\sim$0.1\,ns. This can be
modeled and subtracted using the known libration ephemeris (DE440).
After modeling: residual systematic $\lesssim 0.02$\,ns.

\paragraph{Key point.} Most APOLLO systematics (atmosphere, Earth
orientation, Shapiro delay) are \emph{common-mode} --- they affect both
up-leg and down-leg equally. The nonreciprocal signal is the
\emph{difference}, so common-mode errors cancel.

\subsection{Timeline and Feasibility}

\begin{itemize}
\item \textbf{Near-term (2026--2028)}: CW laser LLR systems (Courde et al.\
  2025) achieve $<$1\,mm precision with higher return rates.
  Still round-trip.
\item \textbf{Artemis era (2028--2032)}: Lunar surface optical clock +
  optical time transfer. One-way measurement becomes feasible.
  The 0.93\,ns signal is 4000$\times$ above the projected noise floor.
\item \textbf{Detection significance}: Single-night detection at $>5\sigma$
  with $\sim$1\,hour of data.
\end{itemize}

%=============================================================================
\section{Finding 3: JILA Atomic Clock LPI Tests}
%=============================================================================

\subsection{What JILA Tests}

Species-dependent gravitational redshift tests compare clocks based on
different atomic species (e.g., Sr vs.\ Yb, or different Sr isotopes)
at the same gravitational potential. The test is parameterized by the
anomaly:
\begin{equation}
\frac{\Delta f}{f} = (1 + \beta_A - \beta_B) \frac{\Delta U}{c^2}
\end{equation}
where $\beta_A - \beta_B$ measures the species-dependent deviation from
the universal gravitational redshift.

Current best limits: $|\beta_{\rm Rb} - \beta_{\rm H}| <
5 \times 10^{-7}$ (Peil et al.\ 2013). Optical clock comparisons aim
for $10^{-8}$--$10^{-9}$.

\subsection{DFD Prediction for LPI Tests}

\paragraph{Derivation chain.}
In DFD, the physical metric is:
\begin{equation}
\tilde{g}_{00} = -c^2 e^{-\psi}
\end{equation}
This gives a \emph{universal} gravitational redshift:
\begin{equation}
\frac{\Delta f}{f} = \frac{\Delta\psi}{2} = \frac{\Delta U}{c^2}
\end{equation}
with $\beta = 0$ for all species at leading order. DFD satisfies the
Weak Equivalence Principle by construction (all test masses fall with
$\vec{a} = (c^2/2)\nabla\psi$).

\paragraph{Species dependence enters through $k_\alpha$.}
The $\psi$-field couples to the fine structure constant via:
\begin{equation}
\alpha(\psi) = \alpha_0(1 + k_\alpha \psi), \qquad
k_\alpha = \frac{\alpha^2}{2\pi} \approx 8.5 \times 10^{-6}
\end{equation}
Different atomic transitions have different sensitivities to $\alpha$
variations. Denoting the sensitivity coefficient $q_A$ for species $A$:
\begin{equation}
\beta_A = 2 k_\alpha q_A
\end{equation}
For typical optical transitions, $q_A \sim 0.01$--$1$, giving:
\begin{equation}
\beta_A \sim 2 \times 8.5 \times 10^{-6} \times q_A
\sim 10^{-5} q_A
\end{equation}

\paragraph{Species-dependent anomaly.}
\begin{equation}
|\beta_A - \beta_B| = 2k_\alpha |q_A - q_B|
\sim 10^{-5} |q_A - q_B|
\end{equation}

For Sr vs.\ Yb: $|q_{\rm Sr} - q_{\rm Yb}| \sim 0.5$--$1$, giving
$|\beta_{\rm Sr} - \beta_{\rm Yb}| \sim 5 \times 10^{-6}$--$10^{-5}$.

\subsection{What This Test Proves and What It Does Not}

\begin{table}[h]
\centering
\begin{tabular}{lcc}
\toprule
\textbf{Aspect} & \textbf{Tested by JILA?} & \textbf{DFD Prediction} \\
\midrule
Universal redshift ($\beta = 0$) & Yes (at $10^{-5}$--$10^{-9}$) &
  $\beta = 0$ at leading order \\
$k_\alpha$ coupling & Yes (at $10^{-5}$ sensitivity) &
  $k_\alpha \approx 8.5 \times 10^{-6}$ \\
$\mu(x)$ crossover & \textbf{No} &
  Requires $a \sim a_0$ environments \\
$\psi$-field existence & \textbf{No} (measures $\alpha$-variation) &
  Indirect constraint only \\
$\lambda$ coupling & \textbf{No} & Requires EM-to-$\psi$ active test \\
Cosmological $\psi_0$ & \textbf{No} & Local test only \\
\bottomrule
\end{tabular}
\end{table}

\begin{tcolorbox}[colback=blue!5!white, colframe=blue!70!black]
\textbf{Precise statement}: A JILA species-dependent LPI test at the
$10^{-5}$ level would be sensitive to the DFD $k_\alpha$ coupling.
A null result at $|\beta_A - \beta_B| < 10^{-6}$ would constrain
$k_\alpha < 5 \times 10^{-7}$ (assuming $|q_A - q_B| \sim 1$), which
would be in $\sim$10$\sigma$ tension with the DFD prediction
$k_\alpha \approx 8.5 \times 10^{-6}$. Detection at $10^{-5}$ would
be \emph{consistent} with but not proof of DFD (other scalar-field
theories predict similar couplings). The test says NOTHING about
$\mu(x)$, $\lambda$, or the cosmological sector.
\end{tcolorbox}

%=============================================================================
\section{Finding 4: H$_0$ Landscape 2025--2026}
%=============================================================================

\subsection{Latest Measurements}

\begin{table}[h]
\centering
\begin{tabular}{lcc}
\toprule
\textbf{Probe} & \textbf{$H_0$ (km/s/Mpc)} & \textbf{Reference} \\
\midrule
SH0ES Cepheids (JWST+HST) & $73.2 \pm 0.9$ & Riess et al.\ 2025 \\
CCHP TRGB (JWST) & $70.4 \pm 1.8$ & Freedman et al.\ 2025 \\
CCHP JAGB (JWST only) & $67.8 \pm 2.7$ & Freedman et al.\ 2025 \\
Planck CMB ($\Lambda$CDM) & $67.4 \pm 0.5$ & Planck 2018 \\
DESI DR2 BAO + CMB & $\sim$68.5 & DESI DR2, March 2025 \\
TDCOSMO lensing & $71.6^{+3.8}_{-4.9}$ & TDCOSMO 2025 \\
\bottomrule
\end{tabular}
\end{table}

\subsection{Assessment}

\begin{itemize}
\item \textbf{Tension persists at $\sim$5$\sigma$}: SH0ES ($73.2$) vs.\
  Planck ($67.4$), $\Delta = 5.8$\,km/s/Mpc.
\item \textbf{JWST ruled out Cepheid crowding}: JWST observations in
  uncrowded host galaxies confirm HST Cepheid distances.
  Crowding/blending is not the source of the tension.
\item \textbf{Freedman CCHP value ($70.4$) is intermediate}: Consistent
  with both SH0ES and Planck at $\sim$1.5$\sigma$ each. Does not resolve
  the tension.
\item \textbf{DESI DR2 + CMB consistent with Planck}: BAO distance
  scale gives $h r_d = 101.54 \pm 0.73$\,Mpc. With standard $r_d$,
  $H_0 \sim 68.5$.
\item \textbf{DESI hints at dynamical dark energy}: 2.8--4.2$\sigma$
  preference for $w(z) \neq -1$ depending on SNe dataset.
  In DFD language, this is the $\psi$-screen redshift dependence
  $\Delta\psi(z)$ being fit to a CPL parameterization.
\end{itemize}

\paragraph{DFD implications.}
DFD predicts $H_0^{\rm TD} < H_0^{\rm SN}$ by 1--2\,km/s/Mpc. Current
TDCOSMO: $71.6^{+3.8}_{-4.9}$. SH0ES: $73.2 \pm 0.9$.
$\Delta = 1.6 \pm 4.0$\,km/s/Mpc --- consistent with DFD prediction
but not yet discriminating. Need TDCOSMO precision $<$2\,km/s/Mpc.

%=============================================================================
\section{Finding 5: BAO Sound Horizon with $c_s = c$}
%=============================================================================

\subsection{The Question}

With the dust-branch confirmation $c_s = c$ for the $\psi$-scalar sector
(W4i17--i18), does the $\psi$-field modify the BAO sound horizon $r_s$?

\subsection{Derivation}

\paragraph{Step 1: What determines $r_s$.}
The BAO sound horizon is the comoving distance traveled by sound waves
in the baryon-photon fluid from the Big Bang to the drag epoch ($z_d
\approx 1060$):
\begin{equation}
r_s = \int_0^{t_d} \frac{c_s^{\rm b\gamma}}{a(t)}\,dt
= \int_{z_d}^{\infty} \frac{c_s^{\rm b\gamma}}{H(z)}\,dz
\end{equation}
where $c_s^{\rm b\gamma} = c/\sqrt{3(1 + R_b)}$ is the baryon-photon
sound speed, with $R_b = 3\rho_b/(4\rho_\gamma)$.

\paragraph{Step 2: Does $\psi$ modify $c_s^{\rm b\gamma}$?}
In DFD, the baryon-photon sound speed depends on:
\begin{itemize}
\item The electromagnetic propagation speed $c_{\rm EM} = c\,e^{-\psi}$
  in the local $\psi$-field.
\item The baryon loading ratio $R_b$, set by BBN.
\end{itemize}
In the homogeneous universe, $\nabla\psi = 0$ (DFD does not modify the
Friedmann equation, W4i14). The background $\psi_0(t)$ enters only
through a uniform rescaling of $c_{\rm EM}$, which is absorbed into the
definition of conformal time. Therefore:

\begin{equation}
\boxed{r_s^{\rm DFD} = r_s^{\Lambda\text{CDM}} + \mathcal{O}(\delta\psi)}
\end{equation}

The leading-order sound horizon is \emph{identical} to $\Lambda$CDM.

\paragraph{Step 3: What about $\psi$-perturbations?}
The $c_s = c$ result means $\psi$-perturbations propagate at $c$, not
as pressureless dust. This means:
\begin{itemize}
\item $\psi$-perturbations do \emph{not} cluster with baryons during
  the acoustic phase.
\item The $\psi$-field provides no additional gravitational potential
  wells for the baryon-photon fluid.
\item BAO oscillations are driven by baryon self-gravity + photon
  pressure, exactly as in $\Lambda$CDM.
\end{itemize}

\paragraph{Step 4: Quantitative modification.}
The $\psi$-field modifies BAO only through $G_{\rm eff}$:
\begin{equation}
G_{\rm eff}(k) = G_N \frac{k^2 + k_\mu^2}{k^2}
\end{equation}
At BAO scales ($k \sim 0.01$--$0.1\;h/$Mpc), $k \gg k_\mu$
(where $k_\mu = a_\star/c \sim 10^{-4}\;h/$Mpc), so:
\begin{equation}
\frac{G_{\rm eff}}{G_N} - 1 = \frac{k_\mu^2}{k^2}
\sim 10^{-6}\text{--}10^{-4}
\end{equation}

\begin{tcolorbox}[colback=green!5!white, colframe=green!70!black]
\textbf{Result}: DFD modifies the BAO sound horizon by
$\Delta r_s / r_s \lesssim 10^{-4}$, far below the DESI DR2 measurement
precision of 0.65\% (isotropic BAO at $z_{\rm eff} = 2.33$).
DFD is indistinguishable from $\Lambda$CDM for all BAO observables.
\end{tcolorbox}

\subsection{Implications for DESI DR2}

\begin{enumerate}
\item \textbf{$D_H/r_d$ and $D_M/r_d$}: Unmodified. DFD predicts the
  same geometric distances as $\Lambda$CDM because it does not modify
  $H(z)$ or $r_d$.

\item \textbf{Dynamical dark energy hints}: DESI's 2.8--4.2$\sigma$
  preference for $w(z) \neq -1$ arises from SNe data combined with BAO.
  In DFD, the SNe distances carry $\psi$-screen bias; BAO distances do
  not. When a CPL model is fit to the combination, the $\psi$-screen
  bias on SNe mimics $w_0 \neq -1$, $w_a \neq 0$. DFD predicts
  $w_0 \sim -0.97$, $w_a \sim +0.06$ (W4i14).

\item \textbf{The $h r_d$ degeneracy}: DESI alone measures $h r_d =
  101.54 \pm 0.73$\,Mpc. DFD does not break this degeneracy differently
  from $\Lambda$CDM. The $\psi$-screen affects only $D_L^{\rm EM}$
  (SNe), not $D_A^{\rm BAO}$ or $D_H^{\rm BAO}$.
\end{enumerate}

%=============================================================================
\section{Synthesis: P5 Status After i19}
%=============================================================================

\begin{table}[h]
\centering
\begin{tabular}{lccc}
\toprule
\textbf{P5 Prediction} & \textbf{Status} & \textbf{Cosmo?} & \textbf{Changed?} \\
\midrule
Lunar NR 0.93\,ns & ALIVE (SNR 4000) & No & Protocol detailed \\
Bubble timing 4.0\,ns & ALIVE & No & No change \\
Pioneer anomaly $H_0 c$ & ALIVE (ungrounded) & No & No change \\
Earth-Mars ranging & ALIVE & No & No change \\
Sidereal discriminant & ALIVE & No & No change \\
$H_0$ tension (7\%) & HONEST NEGATIVE & Yes & \textbf{Confirmed limited} \\
$H_0^{\rm TD} < H_0^{\rm SN}$ & NEW PREDICTION & No$^*$ & NEW (testable now) \\
BAO $r_s$ unmodified & CONFIRMED & Yes & \textbf{NEW derivation} \\
$k_\alpha$ via JILA clocks & TESTABLE & No & \textbf{NEW analysis} \\
\bottomrule
\end{tabular}
\end{table}

$^*$The $H_0^{\rm TD} < H_0^{\rm SN}$ prediction depends on the
$\psi$-screen being real, but is testable with local data and does not
require cosmological $\psi$ clustering.

%=============================================================================
\section{Corrections to Previous Iterations}
%=============================================================================

\begin{enumerate}
\item \textbf{i18 SNR $\sim 4000$ confirmed}: The systematic error
  budget now demonstrates this is robust. Dominant systematic
  (libration) is modelable to $<$0.02\,ns residual.

\item \textbf{i17 $\Delta H_0 \sim 0.4$\,km/s/Mpc confirmed}: No
  secondary mechanism found. This is now FINAL.

\item \textbf{i14 distance-bias framework REINFORCED}: BAO $r_s$
  derivation confirms DFD is a distance-bias theory. BAO unaffected.
  SNe biased. This is the cleanest statement of DFD cosmology.

\item \textbf{Previous $r_s$ discussion}: Appendix J of v3.2 defines
  $r_s = \int c_s(\psi)\,d\tau$ where $c_s(\psi) = c(\psi)/\sqrt{3}$.
  With homogeneous $\psi_0$, $c(\psi_0) = c_0 e^{-\psi_0}$ is spatially
  uniform, absorbed into conformal time. No modification to $r_s/r_d$
  ratios measured by BAO. Consistent with W4i14.
\end{enumerate}

%=============================================================================
\section{Open Questions for i20}
%=============================================================================

\begin{enumerate}
\item \textbf{JILA $k_\alpha$ timeline}: When will Sr--Yb comparisons
  reach $10^{-6}$ sensitivity to $|\beta_{\rm Sr} - \beta_{\rm Yb}|$?
  This would be the first test sensitive to DFD's $k_\alpha$ prediction.
  The FOCOS space mission (proposed) would reach $10^{-9}$.

\item \textbf{Artemis one-way timing}: What is the realistic timeline
  for a lunar optical clock? ESA Moonlight programme? NASA Artemis III+?
  This determines when LLR 0.93\,ns becomes measurable.

\item \textbf{TDCOSMO precision}: When will $H_0^{\rm TD}$ reach
  $<$2\,km/s/Mpc precision to test $H_0^{\rm TD} < H_0^{\rm SN}$?

\item \textbf{mochi\_class urgency remains}: The BAO derivation confirms
  DFD matches $\Lambda$CDM at leading order, but full numerical
  verification requires the DFD Boltzmann code. This is NOT a new
  finding --- it reinforces the i17 urgency.
\end{enumerate}

%=============================================================================
\section*{FINAL ASSESSMENT}
%=============================================================================

\begin{tcolorbox}[colback=blue!5!white, colframe=blue!70!black]
\textbf{P5 remains ALIVE with reduced cosmological ambitions.}

Core strengths: LLR 0.93\,ns (SNR 4000, crisis-immune), BAO survival
(DFD indistinguishable from $\Lambda$CDM), JILA $k_\alpha$ test
(new falsifiable channel).

Core weakness: H$_0$ tension contribution is only 7\%. DFD does not
solve the Hubble tension. This must be honestly communicated.

Net: P5's value proposition rests on LOCAL timing predictions
(LLR, bubble, Pioneer), not cosmological claims. The cosmological sector
survives scrutiny (BAO unmodified, DESI-consistent) but does not
distinguish DFD from $\Lambda$CDM at current precision.
\end{tcolorbox}

\end{document}
